\documentclass{article}
\usepackage[a4paper, margin=1in]{geometry} % Adjust margins
\usepackage{amsmath, amssymb, amsthm}
\usepackage{tikz}
\usepackage{xcolor}
\usepackage{graphicx}
\usepackage{listings}
\usepackage{fancyhdr}
\usepackage{booktabs}
\usepackage{caption}
\usetikzlibrary{automata, positioning}
\usepackage{helvet} % Match font to reference
\pagestyle{fancy}
\usepackage{enumitem,xcolor}


% Ensure proper header spacing
\setlength{\headheight}{15pt} % Adjust height of the header
\setlength{\headsep}{10pt} 

% Define Header/Footer
\fancyhf{}
\fancyhead[L]{\textbf{Advanced Topics in Computer Science}} % Left header
\fancyhead[R]{DFA Minimization} % Right header
\cfoot{\thepage} % Page number at center bottom

% Define colors for section and subsection titles
\usepackage{titlesec}
\definecolor{sectioncolor}{rgb}{0.1, 0.3, 0.7} % Dark blue
\titleformat{\section}{\large\bfseries\color{sectioncolor}}{\thesection}{1em}{}
\titleformat{\subsection}{\bfseries\color{sectioncolor}}{\thesubsection}{1em}{}

% Define custom colors for different text types
\usepackage{tcolorbox}
\tcbuselibrary{theorems}
\newcommand{\important}[1]{\textcolor{red}{#1}}
\newtcbtheorem
  []% init options
  {definition}% name
  {Definition}% title
  {%
    colback=blue!5,
    colframe=blue!35!black,
    fonttitle=\bfseries,
  }% options
  {def}% prefix
\newcommand{\note}[1]{\textcolor{green!50!black}{#1}}

% Theorem environment setup
\newtheorem{theorem}{Theorem}
\newtheorem{lemma}{Lemma}
\newtheorem{corollary}{Corollary}

\definecolor{codegray}{gray}{0.9}
\lstdefinestyle{customc}{
  backgroundcolor=\color{codegray},
  language=C,
  basicstyle=\ttfamily\footnotesize,
  keywordstyle=\color{blue},
  commentstyle=\color{green!50!black},
  stringstyle=\color{red}
}
\begin{document}

\title{\textbf{DFA Minimizer}}
\author{Doodlebob}
\date{}
\maketitle
\thispagestyle{fancy}

The following is an auto-explainer for DFA Minimization as represented by a series of figures depicting the evolving partitions of the DFA’s state set. \\

A Deterministic Finite Automaton (DFA) is denoted by the quintuple
\[
M = (Q, \Sigma, \delta, q_0, F)
\]
and the minimization process works by iteratively refining partitions of the state set until all equivalent states are grouped together. In the minimization process:
\begin{itemize}
    \item The initial partition divides $Q$ into two blocks: the final states $F$ and the non-final states $Q \setminus F$.
    \item At each subsequent step, the partition is refined by examining the transitions for every input symbol and grouping together states that behave equivalently.
    \item This iterative refinement continues until no further splitting is possible, thus yielding a minimal DFA.
\end{itemize}

\section{DFA Definition and Minimization Overview}

A Deterministic Finite Automaton is defined as a 5-tuple 
\[
(Q, \Sigma, \delta, q_0, F)
\]
where:
\begin{itemize}
    \item $Q$ is a finite set of states,
    \item $\Sigma$ is a finite input alphabet,
    \item $\delta : Q \times \Sigma \rightarrow Q$ is the transition function,
    \item $q_0 \in Q$ is the start state, and
    \item $F \subseteq Q$ is the set of final (accepting) states.
\end{itemize}

The minimization of a DFA entails computing an equivalent DFA with the fewest possible states. The process begins by partitioning the state set into final and non-final states and then repeatedly refining these partitions by identifying states that are indistinguishable with respect to the input symbols. The resulting minimized DFA retains the same language as the original.
\section{DFA Transition Function}

The DFA's transition function, $\delta$, is defined as follows:
\begin{align*}\delta(1,\,a) &= 2\\[0.5em]\delta(1,\,b) &= 4\\[1em]\delta(2,\,a) &= 4\\[0.5em]\delta(2,\,b) &= 3\\[1em]\delta(3,\,a) &= 3\\[0.5em]\delta(3,\,b) &= 3\\[1em]\delta(4,\,a) &= 4\\[0.5em]\delta(4,\,b) &= 5\\[1em]\delta(5,\,a) &= 5\\[0.5em]\delta(5,\,b) &= 5\\[1em]\delta(6,\,a) &= 6\\[0.5em]\delta(6,\,b) &= 5\end{align*}
The set of final states is $F = \{ 3, 5 \}$.

\section{Table-Filling (Indistinguishability) Algorithm}

We apply the table-filling algorithm with systematic pair-by-pair analysis. The algorithm begins with an initialization phase that identifies all pairs differing in finality, followed by iterative refinement phases that identify additional distinguishable pairs.
\subsection*{Step 1: Initialization Phase}
We systematically examine all 15 state pairs to identify those that differ in finality. The final states are F = {3, 5}. For each pair (p, q), we check whether exactly one of p or q is final. Our analysis proceeds as follows: (1, 2): 1 is non-final, 2 is non-final; (1, 3): 1 is non-final, 3 is final-> MARK; (1, 4): 1 is non-final, 4 is non-final; (1, 5): 1 is non-final, 5 is final-> MARK; (1, 6): 1 is non-final, 6 is non-final; (2, 3): 2 is non-final, 3 is final-> MARK; (2, 4): 2 is non-final, 4 is non-final; (2, 5): 2 is non-final, 5 is final-> MARK; (2, 6): 2 is non-final, 6 is non-final; (3, 4): 3 is final, 4 is non-final-> MARK; (3, 5): 3 is final, 5 is final; (3, 6): 3 is final, 6 is non-final-> MARK; (4, 5): 4 is non-final, 5 is final-> MARK; (4, 6): 4 is non-final, 6 is non-final; (5, 6): 5 is final, 6 is non-final-> MARK. Therefore, we mark 8 pairs as distinguishable: (1, 3), (1, 5), (2, 3), (2, 5), (3, 4), (3, 6), (4, 5), (5, 6). These receive X marks in our indistinguishability table.
\\[0.5em]
\textbf{Detailed Analysis:}
\begin{itemize}
  \item $(1, 2): 1 is non-final, 2 is non-final$
  \item $(1, 3): 1 is non-final, 3 is final$ $\rightarrow$ \textbf{MARK}
  \item $(1, 4): 1 is non-final, 4 is non-final$
  \item $(1, 5): 1 is non-final, 5 is final$ $\rightarrow$ \textbf{MARK}
  \item $(1, 6): 1 is non-final, 6 is non-final$
  \item $(2, 3): 2 is non-final, 3 is final$ $\rightarrow$ \textbf{MARK}
  \item $(2, 4): 2 is non-final, 4 is non-final$
  \item $(2, 5): 2 is non-final, 5 is final$ $\rightarrow$ \textbf{MARK}
  \item $(2, 6): 2 is non-final, 6 is non-final$
  \item $(3, 4): 3 is final, 4 is non-final$ $\rightarrow$ \textbf{MARK}
  \item $(3, 5): 3 is final, 5 is final$
  \item $(3, 6): 3 is final, 6 is non-final$ $\rightarrow$ \textbf{MARK}
  \item $(4, 5): 4 is non-final, 5 is final$ $\rightarrow$ \textbf{MARK}
  \item $(4, 6): 4 is non-final, 6 is non-final$
  \item $(5, 6): 5 is final, 6 is non-final$ $\rightarrow$ \textbf{MARK}
\end{itemize}

\\[0.25em]
\textbf{Metrics:}
\begin{itemize}
  \item Newly marked this step: 8
  \item Cumulative marked: 8 of 15 (53.3\%)
\end{itemize}


\begin{figure}[h!]
\centering
\begin{minipage}{0.48\linewidth}
    \centering
    \includegraphics[width=\linewidth]{plots/indis_step1.png}
    \caption{Step 1: Indistinguishability table showing finality-based marks}
    \label{fig:indis_table_step1}
\end{minipage}\hfill
\begin{minipage}{0.48\linewidth}
    \centering
    \includegraphics[width=\linewidth]{plots/blocks_step1.png}
    \caption{Partition after Step 1 (states colored by equivalence class)}
    \label{fig:partition_step1}
\end{minipage}
\end{figure}


\paragraph{Equivalence Partition after Step 1}
The current partition of states into equivalence classes is shown in Figure~\ref{fig:partition_step1}:
\[
\left\{\{ 3, 5 \}, \{ 2, 1, 4, 6 \}\right\}
\]

\clearpage
\subsection*{Step 2: Refinement Phase}

\textbf{Summary:} Pairs marked this round because their transitions reach marked witnesses.


\paragraph{Part 2.a \;–\; Pair $( 1, 2 )$}

\emph{Analysis.} Found witness on input b: delta(1,b)=4, delta(2,b)=3. Since ( 3, 4 ) is already marked, we mark ( 1, 2 ).

\noindent\textbf{Transition Analysis:}
\begin{center}
\begin{tabular}{c c c c c c}
\hline
$a$ & $\delta(1,a)$ & $\delta(2,a)$ & Witness Pair & Witness Status & Decision \\
\hline
a & 2 & 4 &
(2, 4) &
Unmarked &
Continue \\
b & 4 & 3 &
(3, 4) &
Marked &
Mark \\\hline
\hline
\end{tabular}
\end{center}


\begin{figure}[h!]
\centering
\includegraphics[width=0.85\linewidth]{plots/pair_steps/step2_a_probe2.png}
\caption{Decisive Probe: $\delta(1, b) = 4$, $\delta(2, b) = 3$ 
$\to$ Witness $(3, 4)$ is marked}
\label{fig:step2_a_probe2}
\end{figure}

\noindent\textbf{Final Decision for $( 1, 2 )$:} Mark.
\textbf{Action:} Add ``X'' to the indistinguishability table.

\vspace{0.75em}
\paragraph{Part 2.b \;–\; Pair $( 1, 3 )$}

\emph{Analysis.} Pair ( 1, 3 ) was already marked before this step; no further probing is needed.

\noindent\textbf{Final Decision for $( 1, 3 )$:} Already marked.

\vspace{0.75em}
\paragraph{Part 2.c \;–\; Pair $( 1, 4 )$}

\emph{Analysis.} Found witness on input b: delta(1,b)=4, delta(4,b)=5. Since ( 4, 5 ) is already marked, we mark ( 1, 4 ).

\noindent\textbf{Transition Analysis:}
\begin{center}
\begin{tabular}{c c c c c c}
\hline
$a$ & $\delta(1,a)$ & $\delta(4,a)$ & Witness Pair & Witness Status & Decision \\
\hline
a & 2 & 4 &
(2, 4) &
Unmarked &
Continue \\
b & 4 & 5 &
(4, 5) &
Marked &
Mark \\\hline
\hline
\end{tabular}
\end{center}


\begin{figure}[h!]
\centering
\includegraphics[width=0.85\linewidth]{plots/pair_steps/step2_c_probe2.png}
\caption{Decisive Probe: $\delta(1, b) = 4$, $\delta(4, b) = 5$ 
$\to$ Witness $(4, 5)$ is marked}
\label{fig:step2_c_probe2}
\end{figure}

\noindent\textbf{Final Decision for $( 1, 4 )$:} Mark.
\textbf{Action:} Add ``X'' to the indistinguishability table.

\vspace{0.75em}
\paragraph{Part 2.d \;–\; Pair $( 1, 5 )$}

\emph{Analysis.} Pair ( 1, 5 ) was already marked before this step; no further probing is needed.

\noindent\textbf{Final Decision for $( 1, 5 )$:} Already marked.

\vspace{0.75em}
\paragraph{Part 2.e \;–\; Pair $( 1, 6 )$}

\emph{Analysis.} Found witness on input b: delta(1,b)=4, delta(6,b)=5. Since ( 4, 5 ) is already marked, we mark ( 1, 6 ).

\noindent\textbf{Transition Analysis:}
\begin{center}
\begin{tabular}{c c c c c c}
\hline
$a$ & $\delta(1,a)$ & $\delta(6,a)$ & Witness Pair & Witness Status & Decision \\
\hline
a & 2 & 6 &
(2, 6) &
Unmarked &
Continue \\
b & 4 & 5 &
(4, 5) &
Marked &
Mark \\\hline
\hline
\end{tabular}
\end{center}


\begin{figure}[h!]
\centering
\includegraphics[width=0.85\linewidth]{plots/pair_steps/step2_e_probe2.png}
\caption{Decisive Probe: $\delta(1, b) = 4$, $\delta(6, b) = 5$ 
$\to$ Witness $(4, 5)$ is marked}
\label{fig:step2_e_probe2}
\end{figure}

\noindent\textbf{Final Decision for $( 1, 6 )$:} Mark.
\textbf{Action:} Add ``X'' to the indistinguishability table.

\vspace{0.75em}
\paragraph{Part 2.f \;–\; Pair $( 2, 3 )$}

\emph{Analysis.} Pair ( 2, 3 ) was already marked before this step; no further probing is needed.

\noindent\textbf{Final Decision for $( 2, 3 )$:} Already marked.

\vspace{0.75em}
\paragraph{Part 2.g \;–\; Pair $( 2, 4 )$}

\emph{Analysis.} Analyze pair ( 2, 4 ) by probing input symbols. No probe led to a marked witness pair; remain unmarked this round.

\noindent\textbf{Transition Analysis:}
\begin{center}
\begin{tabular}{c c c c c c}
\hline
$a$ & $\delta(2,a)$ & $\delta(4,a)$ & Witness Pair & Witness Status & Decision \\
\hline
a & 4 & 4 &
(4, 4) &
Unmarked &
Continue \\
b & 3 & 5 &
(3, 5) &
Unmarked &
Continue \\
\hline
\end{tabular}
\end{center}



\noindent\textbf{Final Decision for $( 2, 4 )$:} Keep unmarked.

\vspace{0.75em}
\paragraph{Part 2.h \;–\; Pair $( 2, 5 )$}

\emph{Analysis.} Pair ( 2, 5 ) was already marked before this step; no further probing is needed.

\noindent\textbf{Final Decision for $( 2, 5 )$:} Already marked.

\vspace{0.75em}
\paragraph{Part 2.i \;–\; Pair $( 2, 6 )$}

\emph{Analysis.} Analyze pair ( 2, 6 ) by probing input symbols. No probe led to a marked witness pair; remain unmarked this round.

\noindent\textbf{Transition Analysis:}
\begin{center}
\begin{tabular}{c c c c c c}
\hline
$a$ & $\delta(2,a)$ & $\delta(6,a)$ & Witness Pair & Witness Status & Decision \\
\hline
a & 4 & 6 &
(4, 6) &
Unmarked &
Continue \\
b & 3 & 5 &
(3, 5) &
Unmarked &
Continue \\
\hline
\end{tabular}
\end{center}



\noindent\textbf{Final Decision for $( 2, 6 )$:} Keep unmarked.

\vspace{0.75em}
\paragraph{Part 2.j \;–\; Pair $( 3, 4 )$}

\emph{Analysis.} Pair ( 3, 4 ) was already marked before this step; no further probing is needed.

\noindent\textbf{Final Decision for $( 3, 4 )$:} Already marked.

\vspace{0.75em}
\paragraph{Part 2.k \;–\; Pair $( 3, 5 )$}

\emph{Analysis.} Analyze pair ( 3, 5 ) by probing input symbols. No probe led to a marked witness pair; remain unmarked this round.

\noindent\textbf{Transition Analysis:}
\begin{center}
\begin{tabular}{c c c c c c}
\hline
$a$ & $\delta(3,a)$ & $\delta(5,a)$ & Witness Pair & Witness Status & Decision \\
\hline
a & 3 & 5 &
(3, 5) &
Unmarked &
Continue \\
b & 3 & 5 &
(3, 5) &
Unmarked &
Continue \\
\hline
\end{tabular}
\end{center}



\noindent\textbf{Final Decision for $( 3, 5 )$:} Keep unmarked.

\vspace{0.75em}
\paragraph{Part 2.l \;–\; Pair $( 3, 6 )$}

\emph{Analysis.} Pair ( 3, 6 ) was already marked before this step; no further probing is needed.

\noindent\textbf{Final Decision for $( 3, 6 )$:} Already marked.

\vspace{0.75em}
\paragraph{Part 2.m \;–\; Pair $( 4, 5 )$}

\emph{Analysis.} Pair ( 4, 5 ) was already marked before this step; no further probing is needed.

\noindent\textbf{Final Decision for $( 4, 5 )$:} Already marked.

\vspace{0.75em}
\paragraph{Part 2.n \;–\; Pair $( 4, 6 )$}

\emph{Analysis.} Analyze pair ( 4, 6 ) by probing input symbols. No probe led to a marked witness pair; remain unmarked this round.

\noindent\textbf{Transition Analysis:}
\begin{center}
\begin{tabular}{c c c c c c}
\hline
$a$ & $\delta(4,a)$ & $\delta(6,a)$ & Witness Pair & Witness Status & Decision \\
\hline
a & 4 & 6 &
(4, 6) &
Unmarked &
Continue \\
b & 5 & 5 &
(5, 5) &
Unmarked &
Continue \\
\hline
\end{tabular}
\end{center}



\noindent\textbf{Final Decision for $( 4, 6 )$:} Keep unmarked.

\vspace{0.75em}
\paragraph{Part 2.o \;–\; Pair $( 5, 6 )$}

\emph{Analysis.} Pair ( 5, 6 ) was already marked before this step; no further probing is needed.

\noindent\textbf{Final Decision for $( 5, 6 )$:} Already marked.

\vspace{0.75em}

\\[0.25em]
\textbf{Metrics:}
\begin{itemize}
  \item Newly marked this step: 3
  \item Cumulative marked: 11 of 15 (73.3\%)
\end{itemize}


\begin{figure}[h!]
\centering
\begin{minipage}{0.48\linewidth}
    \centering
    \includegraphics[width=\linewidth]{plots/indis_step2.png}
    \caption{Step 2: Indistinguishability table (bright red ``X'' = newly marked, dark red ``X'' = previously marked)}
    \label{fig:indis_table_step2}
\end{minipage}\hfill
\begin{minipage}{0.48\linewidth}
    \centering
    \includegraphics[width=\linewidth]{plots/blocks_step2.png}
    \caption{Partition after Step 2 (states colored by equivalence class)}
    \label{fig:partition_step2}
\end{minipage}
\end{figure}


\paragraph{Equivalence Partition after Step 2}
The current partition of states into equivalence classes is shown in Figure~\ref{fig:partition_step2}:
\[
\left\{\{ 1 \}, \{ 3, 5 \}, \{ 2, 4, 6 \}\right\}
\]

\clearpage
\subsection*{Step 3: Refinement Phase}

\textbf{Summary:} No new pairs marked; reached a fixed point.


\paragraph{Part 3.a \;–\; Pair $( 1, 2 )$}

\emph{Analysis.} Pair ( 1, 2 ) was already marked before this step; no further probing is needed.

\noindent\textbf{Final Decision for $( 1, 2 )$:} Already marked.

\vspace{0.75em}
\paragraph{Part 3.b \;–\; Pair $( 1, 3 )$}

\emph{Analysis.} Pair ( 1, 3 ) was already marked before this step; no further probing is needed.

\noindent\textbf{Final Decision for $( 1, 3 )$:} Already marked.

\vspace{0.75em}
\paragraph{Part 3.c \;–\; Pair $( 1, 4 )$}

\emph{Analysis.} Pair ( 1, 4 ) was already marked before this step; no further probing is needed.

\noindent\textbf{Final Decision for $( 1, 4 )$:} Already marked.

\vspace{0.75em}
\paragraph{Part 3.d \;–\; Pair $( 1, 5 )$}

\emph{Analysis.} Pair ( 1, 5 ) was already marked before this step; no further probing is needed.

\noindent\textbf{Final Decision for $( 1, 5 )$:} Already marked.

\vspace{0.75em}
\paragraph{Part 3.e \;–\; Pair $( 1, 6 )$}

\emph{Analysis.} Pair ( 1, 6 ) was already marked before this step; no further probing is needed.

\noindent\textbf{Final Decision for $( 1, 6 )$:} Already marked.

\vspace{0.75em}
\paragraph{Part 3.f \;–\; Pair $( 2, 3 )$}

\emph{Analysis.} Pair ( 2, 3 ) was already marked before this step; no further probing is needed.

\noindent\textbf{Final Decision for $( 2, 3 )$:} Already marked.

\vspace{0.75em}
\paragraph{Part 3.g \;–\; Pair $( 2, 4 )$}

\emph{Analysis.} Analyze pair ( 2, 4 ) by probing input symbols. No probe led to a marked witness pair; remain unmarked this round.

\noindent\textbf{Transition Analysis:}
\begin{center}
\begin{tabular}{c c c c c c}
\hline
$a$ & $\delta(2,a)$ & $\delta(4,a)$ & Witness Pair & Witness Status & Decision \\
\hline
a & 4 & 4 &
(4, 4) &
Unmarked &
Continue \\
b & 3 & 5 &
(3, 5) &
Unmarked &
Continue \\
\hline
\end{tabular}
\end{center}



\noindent\textbf{Final Decision for $( 2, 4 )$:} Keep unmarked.

\vspace{0.75em}
\paragraph{Part 3.h \;–\; Pair $( 2, 5 )$}

\emph{Analysis.} Pair ( 2, 5 ) was already marked before this step; no further probing is needed.

\noindent\textbf{Final Decision for $( 2, 5 )$:} Already marked.

\vspace{0.75em}
\paragraph{Part 3.i \;–\; Pair $( 2, 6 )$}

\emph{Analysis.} Analyze pair ( 2, 6 ) by probing input symbols. No probe led to a marked witness pair; remain unmarked this round.

\noindent\textbf{Transition Analysis:}
\begin{center}
\begin{tabular}{c c c c c c}
\hline
$a$ & $\delta(2,a)$ & $\delta(6,a)$ & Witness Pair & Witness Status & Decision \\
\hline
a & 4 & 6 &
(4, 6) &
Unmarked &
Continue \\
b & 3 & 5 &
(3, 5) &
Unmarked &
Continue \\
\hline
\end{tabular}
\end{center}



\noindent\textbf{Final Decision for $( 2, 6 )$:} Keep unmarked.

\vspace{0.75em}
\paragraph{Part 3.j \;–\; Pair $( 3, 4 )$}

\emph{Analysis.} Pair ( 3, 4 ) was already marked before this step; no further probing is needed.

\noindent\textbf{Final Decision for $( 3, 4 )$:} Already marked.

\vspace{0.75em}
\paragraph{Part 3.k \;–\; Pair $( 3, 5 )$}

\emph{Analysis.} Analyze pair ( 3, 5 ) by probing input symbols. No probe led to a marked witness pair; remain unmarked this round.

\noindent\textbf{Transition Analysis:}
\begin{center}
\begin{tabular}{c c c c c c}
\hline
$a$ & $\delta(3,a)$ & $\delta(5,a)$ & Witness Pair & Witness Status & Decision \\
\hline
a & 3 & 5 &
(3, 5) &
Unmarked &
Continue \\
b & 3 & 5 &
(3, 5) &
Unmarked &
Continue \\
\hline
\end{tabular}
\end{center}



\noindent\textbf{Final Decision for $( 3, 5 )$:} Keep unmarked.

\vspace{0.75em}
\paragraph{Part 3.l \;–\; Pair $( 3, 6 )$}

\emph{Analysis.} Pair ( 3, 6 ) was already marked before this step; no further probing is needed.

\noindent\textbf{Final Decision for $( 3, 6 )$:} Already marked.

\vspace{0.75em}
\paragraph{Part 3.m \;–\; Pair $( 4, 5 )$}

\emph{Analysis.} Pair ( 4, 5 ) was already marked before this step; no further probing is needed.

\noindent\textbf{Final Decision for $( 4, 5 )$:} Already marked.

\vspace{0.75em}
\paragraph{Part 3.n \;–\; Pair $( 4, 6 )$}

\emph{Analysis.} Analyze pair ( 4, 6 ) by probing input symbols. No probe led to a marked witness pair; remain unmarked this round.

\noindent\textbf{Transition Analysis:}
\begin{center}
\begin{tabular}{c c c c c c}
\hline
$a$ & $\delta(4,a)$ & $\delta(6,a)$ & Witness Pair & Witness Status & Decision \\
\hline
a & 4 & 6 &
(4, 6) &
Unmarked &
Continue \\
b & 5 & 5 &
(5, 5) &
Unmarked &
Continue \\
\hline
\end{tabular}
\end{center}



\noindent\textbf{Final Decision for $( 4, 6 )$:} Keep unmarked.

\vspace{0.75em}
\paragraph{Part 3.o \;–\; Pair $( 5, 6 )$}

\emph{Analysis.} Pair ( 5, 6 ) was already marked before this step; no further probing is needed.

\noindent\textbf{Final Decision for $( 5, 6 )$:} Already marked.

\vspace{0.75em}

\\[0.25em]
\textbf{Metrics:}
\begin{itemize}
  \item Newly marked this step: 0
  \item Cumulative marked: 11 of 15 (73.3\%)
\end{itemize}


\begin{figure}[h!]
\centering
\begin{minipage}{0.48\linewidth}
    \centering
    \includegraphics[width=\linewidth]{plots/indis_step3.png}
    \caption{Step 3: Indistinguishability table (bright red ``X'' = newly marked, dark red ``X'' = previously marked)}
    \label{fig:indis_table_step3}
\end{minipage}\hfill
\begin{minipage}{0.48\linewidth}
    \centering
    \includegraphics[width=\linewidth]{plots/blocks_step3.png}
    \caption{Partition after Step 3 (states colored by equivalence class)}
    \label{fig:partition_step3}
\end{minipage}
\end{figure}


\paragraph{Equivalence Partition after Step 3}
The current partition of states into equivalence classes is shown in Figure~\ref{fig:partition_step3}:
\[
\left\{\{ 1 \}, \{ 3, 5 \}, \{ 2, 4, 6 \}\right\}
\]

\clearpage

\section{Minimized DFA}

After reaching the fixed point, all unmarked pairs represent equivalent states that can be merged. The resulting minimized DFA is shown in Figure~\ref{fig:minimized_dfa}:
\begin{figure}[h!]
\centering
\includegraphics[width=0.8\linewidth]{minimized_dfa.png}
\caption{Minimized DFA - equivalent states have been merged into single representatives}
\label{fig:minimized_dfa}
\end{figure}

\subsection*{Summary}

The table-filling algorithm successfully identified all distinguishable state pairs through systematic analysis:
\begin{itemize}
\item \textbf{Initialization (Step 1):} Marked pairs with different finality status through cumulative analysis as shown in Figure~\ref{fig:indis_table_step1}
\item \textbf{Refinement steps:} Iteratively marked pairs whose transitions lead to already-distinguishable pairs
\item \textbf{Convergence:} Reached fixed point when no new pairs could be marked
\item \textbf{Result:} Minimized DFA with merged equivalent states (Figure~\ref{fig:minimized_dfa})
\end{itemize}

\subsection*{Figure Reference Summary}
The following figures were generated during the analysis:
\begin{itemize}
  \item Step 1 table: Figure~\ref{fig:indis_table_step1}
  \item Step 1 partition: Figure~\ref{fig:partition_step1}
  \item Step 2 table: Figure~\ref{fig:indis_table_step2}
  \item Step 2 partition: Figure~\ref{fig:partition_step2}
  \item Decisive probe 1-2 on b: Figure~\ref{fig:step2_a_probe2}
  \item Decisive probe 1-4 on b: Figure~\ref{fig:step2_c_probe2}
  \item Decisive probe 1-6 on b: Figure~\ref{fig:step2_e_probe2}
  \item Step 3 table: Figure~\ref{fig:indis_table_step3}
  \item Step 3 partition: Figure~\ref{fig:partition_step3}
  \item Minimized DFA: Figure~\ref{fig:minimized_dfa}
\end{itemize}
\end{document}